\chapter{Construction of Exotic, Symplectic Spheres \\ in Convex 4-Manifolds}
In this chapter, we prove the closed analogue of theorem \ref{prince}.
\begin{thm}\label{king} There exist infinitely many convex symplectic $4$-manifolds, each containing a pair of exotic symplectic spheres.
\end{thm}
The ambient $4$-manifold will be a knot trace. We recall the definition and properties for context. 
\section{Knot Traces}
\begin{defn}Let $K$ be a knot in $\s^3$. The $n$-trace of $K$ is the $4$-manifold obtained by attaching an $n$-framed $2$-handle to $\B^4$ along $K$. It is denoted as $X_n(K)$.
\end{defn}
\begin{obs}We observe that $$
\partial X_n(K) = \s^3_n(K)
$$
\end{obs}
\begin{prop}Suppose $n \neq 0$. Then $Y_n = \s^3_n(K)$ is a rational homology sphere. 
\end{prop}
\proof It suffices to show that $H_1(Y_n, \Q) = H_2(Y_n, \Q) = 0$, as $Y_n$ is connected and closed.
We first compute the homology with integer coefficients and then switch to rational coefficients.
Consider the open sets: $$
U = \s^3 \setminus \nu(K)^{\mathrm{o}} \qquad V = \nu(K) \footnote{We actually consider an enlargement of $\nu(K)$ which is homotopy equivalent to $\nu(K)$}
$$

Then $H_1(U, \Z)$ is generated by a meridian $\mu$ and $H_1(V, \Z)$ is generated by a longitude $\lambda$. Next, examine the following fragment of the Mayer Vietoris exact sequence: 
$$
\begin{tikzcd}
                                 & 0 \arrow[r]                                & H_2(Y_n) \arrow[r, "\partial_1"] & {}            \\
H_1(U\cap V) \arrow[r, "\psi_1"] & H_1(U)\oplus H_1(V) \arrow[r, "\varphi_1"] & H_1(Y_n) \arrow[r, "\partial_0"] & H_0(U \cap V)
\end{tikzcd}
$$ $$
\psi_1(\mu) = n \mu \qquad \psi_1(\lambda) = \lambda
$$
Since $Ker( \partial_0) = 0$\footnote{This can be seen by examining the degree $0$ piece of the sequence}, we obtain
$$
H_1(Y_n) \cong Im(\varphi_1) \cong \Z^2/Im(\psi_1)\cong \Z^2/(\Z \oplus n\Z) \cong \Z/n\Z
$$
In addition $$
H_2(Y_n) \cong Im(\partial_1) \cong Ker(\psi_1) = 0 
$$
We conclude using the Universal coefficient Theorem $$
H_1(Y_n, \Q) \cong \Z/n\Z \otimes \Q = 0
$$as desired.
\endproof
Fix $m<<0$. Let $K_m$ be the knot bounding exotic symplectic disks $D, D'$ we used to prove Theorem \ref{prince}. Consider $X_n = X_n(K_m) = \B^4 \cup H_2$. We build a pair of spheres $S, S'$, by capping off $D, D'$ with the core of the $2$-handle $H_2$.
\begin{figure}[H]
    \centering
    \includegraphics[scale=0.25]{figures/knottrace.jpeg}
    \caption{The sphere $S$ inside $X_n$}
    \label{fig:knottrace}
\end{figure}
The proof is divided in four steps: \\
(A) For $n \leq -3$, $X_n$ admits a convex symplectic structure, \\
(B) $S$ and $S'$ are symplectic spheres, \\
(C) $S$ and $S'$ are topologically isotopic, \\
(D) Suppose $2 | n$ \footnote{This requirement will be crucial for our proof} and $n <<0$, then $S, S'$ are not smoothly isotopic. \\
\textbf{Proof of (A)} \\
We use the following theorem due to Gay \cite{gay:2-handles}.
\begin{thm} \label{hymn}Let $(X, \omega)$ be a convex $4$-manifold. Suppose that $X' = X \cup H_2$ is constructed by attaching a $n$-framed $2$-handle along the transverse push-off of a Legendrian\footnote{With respect to the contact structure induced by a Liouville vector field} knot $K\subseteq \partial X$. If $n \leq tb(K)-1$, then there exist a symplectic form $\omega'$, such that $(X', \omega')$ is a weak symplectic filling.
\end{thm}
\begin{rem}$K_m$ has a Legendrian realization with $tb(\mathfrak{k}) = -2$ (Figure \ref{fig:tra}(a)). So the requirement of the previous Theorem is satisfied for $n \leq -3$.
\end{rem}
Thus for such values of $n$, $X_n$ admits a symplectic form $\omega'$, such that $\partial X_n = \s^3_n(K_m)$ is weakly symplectically filled by $X_n$.
Then by Proposition \ref{osha}, up to a small deformation of $\omega'$, $X_n$ admits a convex symplectic structure, as desired. \\ \\
\textbf{Proof of (B)} \\
We need to prove that $D$ and $D'$ can be capped off symplectically. 
First, observe that the core disk inherits an area form from the $2$-handle, and therefore it is symplectic. 
\begin{prop}\label{tool4} Using the notations above, suppose that $S$ is obtained by capping off a symplectic surface $F\subseteq \B^4$ with the core of the $2$-handle attached along $\partial F = K$. If $F$ is cylindrical near its boundary\footnote{In our case, this technical requirement is satisfied, see the construction in Lemma \ref{cob}}, then $S$ is symplectic in $(X_n, \omega')$.
\end{prop}
The latter, can be proved by analyzing the explicit gluing maps, and it is carried out in section 7.2 of \cite{bowden2010}. \\ \\
\textbf{Proof of (C)} \\
Since $D$ and $D'$ are topologically isotopic (relatively to the boundary), then the spheres obtained by capping them off are also topologically isotopic. \\ \\
\textbf{Proof of (D)} \\
We mimic the same strategy we used to prove \ref{prince}. Assume, for the sake of contradiction, that there exists a diffeomorphism $\varphi: X_n \rightarrow X_n$ such that $\varphi (S') = S$.
\begin{lem}\label{tool8} Let $n<<0$ and either $m<<0$ or $m=0$. Then every self diffeomorphism of $\partial X_n= \s^3_n(K_m)$ is isotopic to the identity.
\end{lem}
\proof
For $m<<0$ or $m= 0$, $\s^3 \setminus K_m$ is an asymmetric hyperbolic manifold \ref{app}. Thus, by Lemma \ref{cuspi}, $\s^3_n(K_m)$ (which is the Dehn filling of $\s^3 \setminus K_m$) is also hyperbolic and asymmetric for $n <<0$.
We conclude using the fact that every diffeomorphism of an hyperbolic $3$ manifold is isotopic to an isometry.
\endproof
Using the same notations employed in the proof of \ref{prince}, we recall that $J$ (depicted in Figure \ref{fig:km}) bounds a disk $\mathfrak{D'}\subseteq \estd$. As $J \subseteq \s^3 \setminus K_m$ is unaffected from Dehn surgery, it still bounds a disk $\mathfrak{D'} \subseteq X_n$. Moreover, $\mathfrak{D'}$ is disjoint from $S'$. Since $\varphi(J) = J$, it follows that $$
\varphi(\mathfrak{D'}) = \mathfrak{D}
$$
is a disk bounding $J$ and disjoint from $S$.
\section{Describing the Double Branched Covering of an $n$-Trace}
In this section, we obstruct the existence of $\mathfrak{D}$, using the same techniques we used in sections \ref{buffet} \ref{summon}. 
\begin{lem}\label{warietto} $H_1(X_n \setminus S) \cong \Z/n\Z$
\end{lem}
\proof Let $\nu(S)$ be a tubular neighbourhood of $S$ inside $X_n$. Consider the following fragment of Mayer-Vietoris exact sequence: 
$$
\begin{tikzcd}
H_2(X_n) \arrow[r] & H_1(\nu(S) \setminus S) \arrow[r, "\phi"] & H_1(X_n \setminus S)\oplus H_1(\nu(S)) \arrow[r] & 0
\end{tikzcd}$$
Since $H_1(\nu(S) \setminus S)$ is generated by a meridian $\mu$ and $\phi$ is surjective, $\mu$ also generates $H_1(X_n \setminus S)$.
Notice that the parallel pushoff of $K_m$, which is $\lambda + n\mu$, bounds a pushoff of the $2$-handle's core disk. As $[\lambda] = 0 \in H_1(X_n \setminus S)$, it follows that $n[\mu] = 0 \in H_1(X_n \setminus S)$, as desired.
\endproof
This leads to the following theorem.
\begin{thm}\label{wario}
If $2|n$, then the double branched covering of $X_n$ along $S$ exists and it is unique.
Moreover $$
\Sigma(X_n, S) = \Sigma(\B^4, D) \cup H'_2
$$ with the $2$-handle attached along the lift of $K$ with framing $\frac{n}{2}$.
\end{thm}
\proof
By the previous lemma, $H_1(X_n \setminus S) \cong \Z/n\Z$ has a unique index $2$ subgroup. Thus we may repeat the same argument used in section \ref{covers} to prove existence and uniqueness. Now we describe the handle decomposition of $\Sigma(X_n, S)$.
Since $X_n = \B^4 \cup H_2$, we obtain: $$
\Sigma(X_n, D \cup (\D^2 \times 0) ) = \Sigma(\B^4, D) \cup \Sigma(\D^2\times \D^2, \\\D^2 \times 0) = \Sigma(\B^4, D) \cup (\D^2\times \D^2) 
$$
Let $\tilde{K}$ be an arbitrary lift of $K$. We observe that the attaching region of $H_2$, which is $\nu(K)$, lifts to $\nu(\tilde{K}) \subseteq \partial\Sigma(\B^4, D)$.
Therefore $\Sigma(X_n, S)$ is constructed from $\Sigma(\B^4, D)$ by attaching a $2$-handle along $\tilde{K}$. It remains to determine the attaching framing.
\begin{obs}The branched covering map, restricted to $\nu(\tilde{K)}$, in complex coordinates, is given by: $$
\nu(\tilde{K)} \cong \s^1 \times \D^2 \rightarrow \nu(K) \cong \s^1 \times \D^2
$$ $$
(x, z) \rightarrow (x, z^2)
$$
\end{obs}
Thus, the induced map on the first homology groups of the boundaries is defined as: $$
\lambda \rightarrow \lambda \qquad \qquad [0 \times \partial \D^2] \rightarrow 2[0 \times \partial \D^2]
$$
This proves that the framing coefficient is $\frac{n}{2}$.
\endproof
As a result, a handle diagram for $\Sigma(X_n, S)$ is obtained by attaching a $\frac{n}{2}$-framed $2$-handle, along $\tilde{K}$, to a handle diagram of $\Sigma(\B^4, D)$.
Let $TB(K)$ denote the maximal Thurston-Bennequin number of $K$. 
\begin{obs}If $$
\frac{n}{2} \leq TB(\tilde{K}) -1
$$
there exists a Legendrian realization $\mathfrak{\tilde{k}}$ of $\tilde{K}$ such that $$
\frac{n}{2} = tb(\mathfrak{\tilde{k}})-1 \footnote{It is obtained with a sequence of (negative) stabilizations}
$$
\end{obs}
Thus, for such values of $n$, $\Sigma(X_n, S)$ admits a Stein structure.
The final argument can be repeated verbatim. We attach a $-1$-framed $2$-handle to $\Sigma(X_n, S)$ along $\tilde{J}$, a lift of $J$, to obtain a sphere with self intersection $-1$. Such sphere is forbidden in Stein $4$-manifolds. Hence we get to the desired contradiction. 




